\begin{lemma}[Power of a power in an ordered field for integer exponents] \label{lemma:power_of_a_power_in_an_ordered_field_for_integer_exponents}
    \TODO{AI generated, verify}
    Let $F$ be a \CrefAndHyperrefIfExist{definition:field}{field}, let $a \in F$, and let $m,n \in \mathbb{Z}$ such that all powers below are defined in the sense of \CrefAndHyperrefIfExist{definition:exponentiation_in_a_field}. Then
    $$\left(a^m\right)^n = a^{mn}.$$
    In particular, if $a \neq 0_F$ or both $m,n \geq 0$, the identity holds.
\end{lemma}

\begin{proof}
    \TODO{AI generated, verify}
If $a = 0_F$, then by convention in \CrefAndHyperrefIfExist{definition:exponentiation_in_a_field}, we have $0_F^k = 0_F$ for all $k > 0$ and $0_F^0 = 1_F$. The identity is verified case by case: if both $m,n > 0$, then $(0_F^m)^n = 0_F^n = 0_F = 0_F^{mn}$; if $m > 0, n = 0$, then $(0_F^m)^0 = (0_F)^0 = 1_F = 0_F^0$; the remaining cases are handled similarly.

Assume $a \neq 0_F$. Then $a$ belongs to the multiplicative group $F^\times$, so we may use the group power laws (\CrefAndHyperrefIfExist{lemma:additivity_of_powers_in_a_group}). For any $m,n \in \mathbb{Z}$, the element $(a^m)^n$ is computed in the cyclic subgroup generated by $a$. If both are nonnegative, then for $k \geq 0$,
$$\underbrace{(a^m) \cdot (a^m) \cdots (a^m)}_{n \text{ times}} = a^{m+n+\dots+n} = a^{mn},$$
by repeated application of \CrefAndHyperrefIfExist{lemma:multiplicativity_of_powers_in_a_field}. For negative exponents, the identity follows from the fact that $(a^{-1})^n = (a^n)^{-1}$ (\CrefAndHyperrefIfExist{lemma:inverse_of_power_equals_power_of_inverse_in_a_group}).
\end{proof}
